\section{Finite-horizon tracking without small labels}
\label{app:finite-time-uniform}

The learning-speed closure also tracks arbitrary bounded labels on a
prescribed finite horizon.  The constant in the resulting $q^{-1}$ bound
is independent of width, but the sufficient order depends on width.
The theorem below specifies that dependence and proves continuation of the
closure; it assumes neither a Gram gap nor successful fitting.

For parameter increments use the mobility norm
\begin{equation}
 |v|_n^2=\frac{\norm{v_1}_F^2}{n}
       +\sum_{\ell=2}^L\norm{v_\ell}_F^2
       +\frac{\norm{v_w}_2^2}{n}.
 \label{eq:ft-metric}
\end{equation}
The paper's distance $d_n$, the sum of the normalized block norms,
satisfies
\[
 |\theta-\vartheta|_n\le d_n(\theta,\vartheta)
       \le\sqrt{L+1}\,|\theta-\vartheta|_n.
\]
Write $\norm{u}_m^2=m^{-1}\sum_a u_a^2$ for sample RMS and
$X=\max_a\norm{x_a}_2/\sqrt d$.

\begin{proposition}[Finite horizons and a joint width--order bound]
\label{thm:finite-time-uniform}
\label{prop:trial-finite-time}
Fix finite data inputs, $L\ge2$, and constants $K,A_0<\infty$.
The layer activations satisfy
\[
 \phi_\ell\in C^{1,1}_{\mathrm{loc}}(\R),\qquad
 a_\ell=|\phi_\ell(0)|<\infty,\qquad
 v_\ell=\norm{\phi_\ell'}_\infty<\infty.
\]
At every width $n$, allow any finite initialization and labels satisfying
\begin{equation}
 \max_{\ell\ge2}\norm{W_0^{(\ell)}}_{\mathrm{op}}\le K,\qquad
 \max_a\frac{\norm{W_0^{(1)}x_a/\sqrt d}_2}{\sqrt n}\le A_0,
 \qquad B_0:=\frac{\norm{w_0}_2}{\sqrt n}\le1,
 \qquad Y:=\norm{y}_m\le1.
 \label{eq:ft-initial}
\end{equation}
Let $\theta_D$ be the canonical dense flow and $\widehat\theta_{n,q}$
the original order-$q$ learning-speed closure from the same initialization.
For each $T<\infty$ there are $R_T,C_T,\Lambda_T\ge1$, depending only
on $T$, the fixed data inputs, $L,K,A_0$, and $(a_\ell,v_\ell)_{\ell=1}^L$,
with the following property.  Set
\begin{equation}
 \ell_{n,T}=\max_\ell\operatorname{Lip}
   \bigl(\phi_\ell';[-R_T\sqrt n,R_T\sqrt n]\bigr),\qquad
 s_{n,T}=1+\sqrt n\,\ell_{n,T},\qquad
 H_{n,T}=\exp(\Lambda_Ts_{n,T}).
 \label{eq:ft-modulus}
\end{equation}
For every width $n\ge1$ and every integer
\begin{equation}
 q\ge\max\{B_0^2H_{n,T}^2,\ s_{n,T}H_{n,T}\},
 \label{eq:ft-order}
\end{equation}
the closure exists uniquely on $[0,T]$ and
\begin{equation}
 \sup_{t\le T}d_n(\widehat\theta_{n,q}(t),\theta_D(t))
 \le C_TH_{n,T}\left(\frac{B_0}{q^{3/2}}
                       +\frac{s_{n,T}}{q^2}\right)
 \le\frac{2C_T}{q}.
 \label{eq:ft-tracking}
\end{equation}
Thus $q\ge\lceil e^{2\Lambda_Ts_{n,T}}\rceil$ is sufficient in general;
if $w_0=0$, it suffices that
$q\ge\lceil s_{n,T}e^{\Lambda_Ts_{n,T}}\rceil$.
The bounds $Y,B_0\le1$ can be replaced by any fixed finite bounds,
with corresponding changes in the constants.
\end{proposition}

For globally Lipschitz gates, $\ell_{n,T}\le J$ for a fixed $J$, so a
sufficient order has $\log q\ge C_T(1+\sqrt n)$.
This is a bound on a joint region of $(n,q)$, rather than a
$C_T/q$ estimate for unrestricted widths and orders.  It also imposes no
smallness requirement on labels relative to a Gram margin.  The large
sufficient orders may exceed $n/m$, so this conclusion alone does not
guarantee a reduction in matrix rank or moving state.

\paragraph{The unchanged autonomous system.}
For completeness, let $p_k(u)=P_k(2u-1)$, where $P_k$ is the Legendre
polynomial normalized by $P_k(1)=1$.  For each hidden link and sample the
closure stores $q$ moments $\bar h_{a,k}^{(\ell-1)}$ and
$\bar\delta_{a,k}^{(\ell)}$.  Its clock satisfies $\dot\tau=\rho$,
$\tau(0)=1$, and either moment $M_k$ obeys
\begin{equation}
 \dot M_k=S-\frac{\rho}{\tau}
   \left[kM_k+\sum_{j<k}(2j+1)M_j\right],\qquad 0\le k<q,
 \label{eq:ft-moments}
\end{equation}
with source $S=\rho h_a^{(\ell-1)}$ or $S=r_a\delta_a^{(\ell)}$.
Only the degree-zero forward moments are initially nonzero, equal to
the initial features.  The hidden matrices are
\begin{equation}
 \widehat W^{(\ell)}=W_0^{(\ell)}-
 \frac{2}{mn\tau}\sum_{a,k<q}(2k+1)
      \bar\delta_{a,k}^{(\ell)}
      (\bar h_{a,k}^{(\ell-1)})^\top.
 \label{eq:ft-reconstruction}
\end{equation}
All fields in these equations come from the reconstructed network;
the first layer and readout follow the canonical dense update formulas.
The raw system is locally Lipschitz at fixed $n,q$, including at
$\rho=0$, because it uses $\rho$ and $r_a\delta_a$, and divides only
by $\tau>0$.  If $\rho(0)=0$, every raw velocity vanishes and both
paths are stationary.  Otherwise backward uniqueness prevents a regular
solution from reaching such an equilibrium in finite time.  Thus the
clock change used below is legitimate on compact regular intervals.

The proof first obtains a uniform region from dense loss dissipation and
stops the closure at distance one from the dense path.  Forward-history
energy then controls the defect relative to the residual, before any
backward derivative estimate is used.  The resulting backward regularity
supplies a second projection factor.  An explicit stability estimate and
a strict distance bound finally remove the stop.

\subsection*{Polynomial estimates and the exact defect}

Let $\Pi_q$ be the $L^2(0,\tau)$ orthogonal projection onto polynomials
of degree below $q$, acting on any finite-dimensional Hilbert space
$\mathcal H$.  The following estimates have dimension-independent
constants.
\begin{lemma}[Legendre tails and endpoint bounds]
\label{lem:ft-polynomial}
\begin{align}
 \norm{(I-\Pi_q)g}_{L^2}^2
 &\le\frac{1}{q(q+1)}\int_0^\tau
       \xi(\tau-\xi)\norm{g'(\xi)}_{\mathcal H}^2\dd\xi
 \le\frac{\tau^2\norm{g'}_{L^2}^2}{4q(q+1)},
 \label{eq:ft-tail}\\
 \norm{g(\tau)-(\Pi_qg)(\tau)}_{\mathcal H}
 &\le C\sqrt{\tau/q}\,\norm{g'}_{L^2},
 \qquad
 \norm{(\Pi_qb)(\tau)}_{\mathcal H}
 \le C\sqrt q\,\norm{b}_{L^\infty(\mathcal H)},
 \label{eq:ft-endpoints}\\
 \norm{(I-\Pi_q)\mathbf1_{[1,\tau]}}_{L^2(0,\tau)}
 &\le C\sqrt{\tau/q}\qquad(\tau\ge1).
 \label{eq:ft-step}
\end{align}
The first two estimates concerning $g$ require $g\in H^1$; the estimate
for $b$ requires only boundedness.
\end{lemma}

\begin{proof}
The weighted tail and endpoint estimates are those of
\Cref{lem:at-projection}, with $A=\tau$; its proof applies to the
same Hilbert-valued histories.  To prove the step estimate,
for $q\ge2$ replace the step in \eqref{eq:ft-step} by a linear
ramp of width $\tau/q$ on a side of its jump with that much room.
One side always has length at least $\tau/2$.  The replacement error is
at most $\sqrt{\tau/q}$ and the derivative norm is $\sqrt{q/\tau}$.
Projection contraction and \eqref{eq:ft-tail} prove the claim.
For $q=1$ contraction suffices, and at $\tau=1$ the step vanishes almost
everywhere.
\end{proof}

On the closure's clock use the constant forward prefix on $[0,1]$ and
the zero backward prefix, followed by the actual histories
$h_a^{(\ell)}$ and $b_a^{(\ell)}=(r_a/\rho)\delta_a^{(\ell)}$.
The moment integrals of these histories obey \eqref{eq:ft-moments},
by the polynomial identity
$u p_k'=kp_k+\sum_{j<k}(2j+1)p_j$; uniqueness therefore identifies them
with the raw moments.  A star denotes projection at the current endpoint.
For $D_g=\int_0^\tau\norm{g-\Pi_qg}_2^2\dd\xi$,
\begin{equation}
 \dot D_g=\rho\norm{g-g^*}_2^2,\qquad D_g(0)=0.
 \label{eq:ft-error-energy}
\end{equation}
Indeed, the derivative of the fitted polynomial lies in the polynomial
space, so its pairing with the error is zero.  Alternatively, in a fixed
monomial basis with Gram matrix $G(\tau)$, the projected pairing of
$b,h$ is $M_bG^{-1}M_h^\top$.  Differentiation using
$M_b'=b(\tau)v(\tau)^\top$ and $G'=v(\tau)v(\tau)^\top$ gives
$b(h^*)^\top+b^*h^\top-b^*(h^*)^\top$.  Applying this identity to
\eqref{eq:ft-reconstruction} yields exactly
\begin{equation}
 \dot{\widehat\theta}=F(\widehat\theta)+E,\quad E_1=E_w=0,
 \qquad E_\ell=\frac{2\rho}{mn}\sum_a
      (b_a^{(\ell)}-b_a^{(\ell)*})
      (h_a^{(\ell-1)}-h_a^{(\ell-1)*})^\top.
 \label{eq:ft-defect}
\end{equation}
These differentiation formulas require only locally integrable histories;
the backward prefix jump causes no extra term.  The rank-one identity
$\norm{uv^\top/n}_F=(\norm u_2/\sqrt n)(\norm v_2/\sqrt n)$ and
Cauchy--Schwarz in \eqref{eq:ft-error-energy} give
\begin{equation}
 \int_0^t\norm{E_\ell}_F\dd s
 \le\frac{2}{mn}\sum_a\sqrt{D_{b,\ell,a}(t)D_{h,\ell-1,a}(t)}.
 \label{eq:ft-absolute-defect}
\end{equation}

\subsection*{A uniform comparison region and forward regularity}

Write $\eta=\rho(0)$.  Since $|\phi_\ell(u)|\le a_\ell+v_\ell|u|$,
the recursion $M_1^0=a_1+v_1A_0$,
$M_\ell^0=a_\ell+v_\ell K M_{\ell-1}^0$ gives
$\eta\le Q_0:=1+M_L^0$.  Dense loss dissipation in the mobility norm is
\[
 \frac{\dd}{\dd t}\rho_D^2=-|F(\theta_D)|_n^2,
 \qquad |\theta_D(t)-\theta_0|_n\le\eta\sqrt t.
\]
The second inequality is Cauchy--Schwarz applied to the integrated first
one.  It bounds every dense coordinate on finite intervals at fixed
width, and the locally Lipschitz dense vector field consequently extends
past any proposed finite maximal time.

Stop the closure at its first mobility distance
$|\widehat\theta-\theta_D|_n=1$, or at its
maximal existence time, within $[0,T]$.  Put $D_T=Q_0\sqrt T$.
Up to this stop both paths, and every parameter segment joining them,
belong to the convex region
\begin{equation}
 \begin{aligned}
 \norm{W^{(\ell)}}_{\mathrm{op}}\le D:=K+D_T+1\quad(\ell\ge2),
 &\qquad\frac{\norm w_2}{\sqrt n}\le B:=2+D_T,\\
 \max_a\frac{\norm{W^{(1)}x_a/\sqrt d}_2}{\sqrt n}
       &\le A:=A_0+X(D_T+1).
 \end{aligned}
 \label{eq:ft-region}
\end{equation}
Define
\[
 M_1=a_1+v_1A,\quad M_\ell=a_\ell+v_\ell DM_{\ell-1},\qquad
 \beta_L=v_LB,\quad\beta_\ell=v_\ell D\beta_{\ell+1}.
\]
Throughout this region, the feature and backward RMS norms at each sample
are at most $M_\ell$ and $\beta_\ell$, respectively.  Choose
\[
 R_T=1+\max\{A,DM_1,\ldots,DM_{L-1}\}.
\]
Every preactivation coordinate in the entire comparison region belongs
to $[-R_T\sqrt n,R_T\sqrt n]$, justifying \eqref{eq:ft-modulus} without
using a local modulus to construct its own domain.  Also
$\rho\le\rho_*:=1+BM_L$, and
$1\le\tau\le A_c:=1+T\rho_*$.
All these constants are independent of $n,q$ and of gate moduli.

In estimates for the closure we suppress hats and use the sample/neuron
Hilbert norm
$\norm{g}_{\mathcal H}^2=(mn)^{-1}\sum_a\norm{g_a}_2^2$.
Let $Z_\ell=\int_0^\tau\norm{\partial_\xi h^{(\ell)}}_{\mathcal H}^2
\dd\xi$ and $e_E=\sum_{\ell=2}^L\norm{E_\ell}_F$.
All forward histories belong to $H^1$ on compact pre-stop intervals:
their physical curves are absolutely continuous, $\rho$ is positive
there, and their prefixes join continuously.
The backward array satisfies $\norm{b^{(\ell)}}_{\mathcal H}\le\beta_\ell$.
Endpoint evaluation on degree-below-$q$ polynomials has norm
$q/\sqrt\tau$, since $\sum_{k<q}(2k+1)=q^2$.
Projection contraction therefore bounds its endpoint error by
$(q+1)\beta_\ell$.  Inserting this into \eqref{eq:ft-defect}, and using
\eqref{eq:ft-error-energy} and \eqref{eq:ft-tail}, gives
\begin{equation}
 \int_0^t\rho\norm{E_\ell/\rho}_F^2\dd s
 \le4(q+1)^2\beta_\ell^2
       \norm{(I-\Pi_q)h^{(\ell-1)}}_{L^2(\mathcal H)}^2
 \le2A_c^2\beta_\ell^2Z_{\ell-1}.
 \label{eq:ft-square-defect}
\end{equation}
The first-layer clock derivative has RMS at most
$2v_1X^2\beta_1$.  For later layers differentiate the forward recurrence:
\[
 \partial_\xi h_a^{(\ell)}=\phi_\ell'(z_a^{(\ell)})\od
 \left[(F_\ell/\rho+E_\ell/\rho)h_a^{(\ell-1)}
                    +W^{(\ell)}\partial_\xi h_a^{(\ell-1)}\right].
\]
Here $\norm{F_\ell/\rho}_F\le2\beta_\ell M_{\ell-1}$.
The three-term squared triangle inequality and
\eqref{eq:ft-square-defect} prove inductively $Z_\ell\le Z_\ell^*$,
where, with $S=T\rho_*$,
\begin{align*}
 Z_1^*&=4Sv_1^2X^4\beta_1^2,\\
 Z_\ell^*&=3v_\ell^2\left[
       4S\beta_\ell^2M_{\ell-1}^4+
       (D^2+2A_c^2\beta_\ell^2M_{\ell-1}^2)Z_{\ell-1}^*\right].
\end{align*}
Thus forward derivative energy is bounded independently of width and order.
By \eqref{eq:ft-endpoints}, the forward endpoint error is at most
$C_Tq^{-1/2}$ and the backward endpoint error at most $C_Tq^{1/2}$,
in $\mathcal H$.  Their product in \eqref{eq:ft-defect} proves
\begin{equation}
 e_E\le C_T\rho,\qquad |\dot{\widehat\theta}|_n\le C_T\rho,
 \qquad\max_{\ell,a}
 \frac{\norm{\dot z_a^{(\ell)}}_2+\norm{\dot h_a^{(\ell)}}_2}{\sqrt n}
       \le C_T\rho.
 \label{eq:ft-relative}
\end{equation}
The last estimate follows by forward differentiation through bounded
operators and slopes.  The same calculation for an arbitrary parameter
increment $v$, including
$Df_a[v]=v_w^\top h_a^{(L)}/n+w^\top Dh_a^{(L)}[v]/n$,
gives $\norm{Df[v]}_m\le C_T|v|_n$ throughout \eqref{eq:ft-region}.
Consequently $\norm{\dot r}_m\le C_T\rho$ and
$|\dot\rho|\le C_T\rho$.  For $\eta>0$ this proves
\begin{equation}
 c_T\eta\le\rho(t)\le C_T\eta,\qquad
 \int_0^t\rho\dd s\le C_T\eta,\qquad
 \norm{\dot c}_m\le C_T\quad(c=r/\rho),\qquad
 Z_\ell(t)\le C_T\eta.
 \label{eq:ft-relative-residual}
\end{equation}
The third bound uses $\dot c=(\dot r-c\dot\rho)/\rho$ and
$\norm c_m=1$; the last uses \eqref{eq:ft-relative} and
$\dd\xi=\rho\dd t$.  The positive lower bound is derived relative
to $\eta$ and is not an additional hypothesis.

\subsection*{The second projection factor and stability}

Integrating the readout equation and using
\eqref{eq:ft-relative-residual} refines the backward bounds to
\begin{equation}
 \frac{\norm{w(t)}_2}{\sqrt n}\le B_0+C_T\eta,
 \qquad\max_{\ell,a}\frac{\norm{\delta_a^{(\ell)}(t)}_2}{\sqrt n}
                         \le C_T(B_0+\eta).
 \label{eq:ft-small-residual-carriers}
\end{equation}
Use the full carriers
$k_a^{(L)}=w$ and
$k_a^{(\ell)}=(W^{(\ell+1)})^\top\delta_a^{(\ell+1)}$.
Their RMS norms are bounded on \eqref{eq:ft-region}, so
$\norm{k_a^{(\ell)}}_\infty\le\norm{k_a^{(\ell)}}_2\le C_T\sqrt n$.
A Lipschitz scalar function composed with an absolutely continuous curve
is absolutely continuous, with derivative bounded almost everywhere by
its Lipschitz constant times the curve speed.  This follows directly
from the disjoint-interval criterion for absolute continuity and then
from difference quotients.  Applying it to $\phi_\ell'$ gives
\[
 \left|\frac{\dd}{\dd t}\phi_\ell'(z_{a,i}^{(\ell)})\right|
        \le\ell_{n,T}|\dot z_{a,i}^{(\ell)}|\quad\text{a.e.}
\]
Differentiate
$\delta_a^{(\ell)}=\phi_\ell'(z_a^{(\ell)})\od k_a^{(\ell)}$.
Its gate-variation term has RMS at most
$C_T\sqrt n\,\ell_{n,T}\rho$ by \eqref{eq:ft-relative}.
For $\ell<L$,
\[
 \dot k_a^{(\ell)}=(\dot W^{(\ell+1)})^\top\delta_a^{(\ell+1)}
                 +(W^{(\ell+1)})^\top\dot\delta_a^{(\ell+1)}.
\]
The first term has RMS at most $C_T\rho$ and the second propagates the
next derivative with coefficient $D$.  Starting with the readout
derivative proves, through fixed depth,
\[
 \max_{\ell,a}\frac{\norm{\dot\delta_a^{(\ell)}}_2}{\sqrt n}
       \le C_Ts_{n,T}\rho,\qquad
 \norm{\dot b^{(\ell)}}_{\mathcal H}
       \le C_T(B_0+s_{n,T}\rho).
\]
For the last bound use $b=c\delta$, \eqref{eq:ft-small-residual-carriers},
$\norm{\dot c}_m\le C_T$, and $\eta\le C_T\rho$.
Only the full-carrier RMS-to-maximum conversion introduces width here;
no tail assumption on a dense carrier is used.

The initial backward jump $b_0^{(\ell)}=b^{(\ell)}(1+)$ has
$\mathcal H$ norm at most $C_TB_0$.  Subtract
$b_0^{(\ell)}\mathbf1_{[1,\tau]}$ from its history to obtain a continuous
remainder $\widetilde b^{(\ell)}$.
Changing variables and using \eqref{eq:ft-relative-residual} gives
\[
 \int_0^\tau\norm{\partial_\xi\widetilde b^{(\ell)}}_{\mathcal H}^2\dd\xi
 =\int_0^t\frac{\norm{\dot b^{(\ell)}}_{\mathcal H}^2}{\rho}\dd s
 \le C_T\left(\frac{B_0^2}{\eta}+s_{n,T}^2\eta\right).
\]
If $Q_h,Q_b$ denote the corresponding $L^2(\mathcal H)$ projection
tails, \eqref{eq:ft-tail}, \eqref{eq:ft-step}, and
\eqref{eq:ft-relative-residual} yield
\[
 Q_{h,\ell}\le\frac{C_T\sqrt\eta}{q},\qquad
 Q_{b,\ell}\le C_T\left[
    \frac{B_0}{\sqrt q}
       +\frac{B_0/\sqrt\eta+s_{n,T}\sqrt\eta}{q}\right].
\]
Sample Cauchy--Schwarz in \eqref{eq:ft-absolute-defect} gives
$\int_0^t\norm{E_\ell}_F\dd s\le2Q_{b,\ell}Q_{h,\ell-1}$.
Multiplying the tails before discarding their residual factors proves
\begin{equation}
 \int_0^t e_E\dd s
 \le C_T\left[\frac{B_0\sqrt\eta}{q^{3/2}}
                      +\frac{B_0+s_{n,T}\eta}{q^2}\right]
 \le C_T\left[\frac{B_0}{q^{3/2}}+\frac{s_{n,T}}{q^2}\right].
 \label{eq:ft-source}
\end{equation}
The inverse residual has canceled.  Thus the constants stay uniform as
$\eta\downarrow0$; the exactly zero case was separated before division.
When $w_0=0$, there is no prefix jump and no $q^{-3/2}$ term.

To compare two parameters $\theta,\vartheta$ in
\eqref{eq:ft-region}, forward subtraction gives
\[
 \max_{\ell,a}\frac{\norm{\Delta z_a^{(\ell)}}_2+
                         \norm{\Delta h_a^{(\ell)}}_2}{\sqrt n}
       \le C_T|\theta-\vartheta|_n.
\]
For example
$\Delta z^{(\ell)}=\Delta W^{(\ell)}h^{(\ell-1)}(\vartheta)
+W^{(\ell)}(\theta)\Delta h^{(\ell-1)}$; the two terms are controlled
by the feature RMS and the operator norm, respectively.
In backward subtraction put the full carrier of $\vartheta$ in the
gate-difference term.  Its maximum norm is at most $C_T\sqrt n$, so
that term has RMS at most
$C_T\sqrt n\,\ell_{n,T}|\theta-\vartheta|_n$.
The matrix-difference term is bounded by its Frobenius norm times a
bounded backward RMS, and the next backward difference propagates
through bounded operators and slopes.  Descending from the readout gives
\[
 \max_{\ell,a}\frac{\norm{\Delta\delta_a^{(\ell)}}_2}{\sqrt n}
       \le C_Ts_{n,T}|\theta-\vartheta|_n.
\]
Subtracting the canonical gradient formulas one factor at a time now
proves
\begin{equation}
 |F(\theta)-F(\vartheta)|_n
       \le C_Ts_{n,T}|\theta-\vartheta|_n.
 \label{eq:ft-stability}
\end{equation}
This finite-difference argument uses locally Lipschitz gates and does
not assume an everywhere-defined second activation derivative.
Subtracting the dense equation from \eqref{eq:ft-defect}, and applying
\eqref{eq:ft-source} and \eqref{eq:ft-stability}, gives on the stopped
interval
\begin{equation}
 \sup_{s\le t}d_n(\widehat\theta(s),\theta_D(s))
 \le C_*e^{a_Ts_{n,T}}
          \left[B_0q^{-3/2}+s_{n,T}q^{-2}\right].
 \label{eq:ft-stopped}
\end{equation}
Indeed the mobility error is bounded by its nondecreasing accumulated defect
plus $C_Ts_{n,T}$ times its time integral; replacing the defect by its
terminal value and applying an integrating factor proves this form of
Gronwall's inequality.  Multiplication by $\sqrt{L+1}$ gives the stated
bound for $d_n$, absorbed in $C_*$.  The constants $C_*,a_T$ are
independent of $n,q$.

Choose $C_T=C_*\ge1$ and
$\Lambda_T\ge\max\{1,a_T,\log(4C_*)\}$ in the theorem.
Then $H_{n,T}\ge s_{n,T}$ and $H_{n,T}\ge4C_*$.
Under \eqref{eq:ft-order}, multiplying the right side of
\eqref{eq:ft-stopped} by $q$ bounds each term by $C_*$, and
$q\ge s_{n,T}H_{n,T}\ge4C_*$.  Hence
\[
 \sup_{s\le t}d_n(\widehat\theta(s),\theta_D(s))
          \le2C_*/q\le1/2.
\]
Since $|\widehat\theta-\theta_D|_n\le d_n(\widehat\theta,\theta_D)$,
continuity rules out a first exit at mobility distance one.  A finite maximal raw
endpoint before $T$ is also impossible: the physical parameters stay
within a bounded radius of the dense path at each fixed width,
$1\le\tau\le A_c$, and every raw moment is the integral of its bounded
finite-width history against $|p_k|\le1$.  Thus the full finite raw state
stays in a compact subset of $\tau>0$.  Its locally Lipschitz field is
bounded there, giving a limiting state and local continuation at a
putative endpoint.  This proves \eqref{eq:ft-tracking} on all of $[0,T]$.
Finally, $H_{n,T}\ge s_{n,T}$ and $B_0\le1$ show that the simpler order
conditions stated after \eqref{eq:ft-tracking} imply
\eqref{eq:ft-order}. \qed

\begin{corollary}[Predictions]
\label{cor:finite-time-predictions}
Under \Cref{thm:finite-time-uniform}, training predictions have the same
$C_T/q$ bound.  If additionally
$\norm{W_0^{(1)}}_F/\sqrt n\le A_1$, then for any probability measure
$\mu$ on inputs with finite second moment,
\[
 \left[\int\sup_{t\le T}
  |\widehat f_{n,q}(t,x)-f_{n,D}(t,x)|^2\dd\mu(x)\right]^{1/2}
 \le\frac{C_{T,A_1}}{q}
       \left[\int(1+\norm x_2/\sqrt d)^2\dd\mu(x)\right]^{1/2}.
\]
The same estimate holds uniformly on any fixed bounded test-input set,
with its radius replacing the moment factor.
\end{corollary}
\begin{proof}
Training predictions use the differential bound already proved on
\eqref{eq:ft-region}.  The additional first-layer bound, dense energy and
the distance bound control $\norm{W^{(1)}}_F/\sqrt n$ on both paths and
their joining segments.  Linear activation growth and the bounded
hidden operators imply feature RMS at most
$C_{T,A_1}(1+\norm x_2/\sqrt d)$ on those segments.  Forward
differentiation using bounded slopes, followed by the readout pairing,
therefore gives
$|Df(x)[v]|\le C_{T,A_1}(1+\norm x_2/\sqrt d)|v|_n$.
Integrate along the parameter segment, apply
\eqref{eq:ft-tracking}, take the pointwise supremum in time, and then
take the $L^2(\mu)$ norm.
\end{proof}

\paragraph{Relation to the fixed-width guarantee.}
The fixed-width learning-speed theorem allows arbitrary
$C^{1,1}_{\mathrm{loc}}$ activations and constants depending on the
particular width and initialization.  The present result supplements it
by assuming bounded slopes and uniform initial scales, and displaying
which width dependence can be absorbed into the order.  For unbounded
activations the present proof establishes closure continuation only in
\eqref{eq:ft-order}; outside that region
\eqref{eq:ft-stopped} is only a stopped estimate.  No all-time estimate,
population existence theorem, or width sampling rate follows from this
finite-horizon argument.
